\section{Análisis del Paper}
\label{sec:analisis-extension}

\subsection{Problema y solución}

Waldspurger y Weihl~\cite{Waldspurger1994} identifican que los mecanismos de
planificación basados en prioridades (incluyendo esquemas de prioridad
dinámica como el de Unix) no ofrecen un control \emph{proporcional} y
predecible sobre la asignación de recursos: ajustar prioridades es una
operación ad-hoc, sin una relación cuantitativa clara entre el valor
asignado y la fracción de recurso efectivamente recibida, y las políticas
que sí intentan dar garantías proporcionales (como \textit{fair-share
scheduling} clásico) suelen requerir mecanismos de retroalimentación
complejos o no componen bien cuando los derechos deben transferirse
temporalmente entre entidades (por ejemplo, de un cliente bloqueado a un
servidor que trabaja en su nombre).

La solución propuesta es la \textbf{planificación por lotería}: cada
entidad recibe una cantidad de \textit{boletos} que representa su derecho
sobre el recurso, y cada decisión de asignación se resuelve sorteando un
boleto al azar entre los boletos activos. El mecanismo es simple de
implementar, componible (los boletos se pueden agrupar, transferir o
subdividir sin lógica especial), y --- a diferencia de los esquemas de
prioridad --- da una garantía \emph{probabilística} de proporcionalidad en
vez de una garantía determinista de corto plazo.

\subsection{Proporcionalidad probabilística}

La proporcionalidad de la lotería es un resultado esperado, no una garantía
por despacho: cada sorteo es un ensayo de Bernoulli independiente, y la
fracción de sorteos ganados por una tarea con $t_i$ boletos de un total
activo $T$ converge a $t_i/T$ por la ley de los grandes números conforme
aumenta el número de sorteos observados, con una varianza que decrece con
ese mismo número~\cite{NIST2003Handbook, Illowsky2024}. Es exactamente lo
que exhibe la Fig.~\ref{fig:convergencia}
(Sección~\ref{sec:resultados}): en pocos cientos de despachos el share
observado todavía oscila con amplitud visible; hacia el final de la ventana
de $10^4$ despachos converge de forma estable a la fracción objetivo. Este
carácter probabilístico --- y no una asignación exacta y determinista de
cada intervalo --- es la diferencia conceptual central entre la
planificación por lotería y un planificador de cuotas fijas.

\subsection{Extensión elegida frente a las demás opciones}

El enunciado ofrece cinco extensiones posibles sobre el mecanismo base
(Tabla~\ref{tab:opciones}); el equipo eligió la Opción A
(\textit{compensation tickets}, Sección~\ref{sec:extension}) por las razones
ya expuestas: reutiliza el mecanismo cooperativo/quantum que el proyecto
base debía implementar de todas formas, y su pregunta de evaluación se
responde con el mismo experimento de proporcionalidad que ya exige el
proyecto.

\begin{table}[htbp]
\caption{Las 5 extensiones propuestas por el enunciado}
\label{tab:opciones}
\renewcommand{\arraystretch}{1.25}
\begin{tabular}{p{1.6cm}p{2.7cm}}
\toprule
\textbf{Opción} & \textbf{Requisito mínimo} \\
\midrule
A. \textit{Compensation tickets} (elegida) & Modelar una tarea que cede antes del quantum e inflar temporalmente su valor \\
\midrule
B. Boletos dinámicos & Aplicar cambios de tickets en despachos definidos por un archivo de eventos \\
\midrule
C. \textit{Ticket transfer} & Modelar cliente bloqueado y servidor que recibe temporalmente sus boletos \\
\midrule
D. \textit{Currencies} & Dos monedas respaldadas por shares base e inflación local \\
\midrule
E. Selección $O(\log n)$ & Árbol de sumas parciales en vez de lista acumulativa \\
\bottomrule
\end{tabular}
\end{table}

Las opciones descartadas se evaluaron y se prefirió A sobre ellas por menor
superficie de bugs de concurrencia nuevos: C requiere modelar estados
cliente/servidor adicionales a los tres estados de \texttt{Task} ya
definidos; D requiere una doble contabilidad de monedas con su propia
lógica de conversión; B modifica la asignación de boletos desde una fuente
externa durante la corrida, lo que interactúa con la validación de entrada
de formas que el enunciado no especifica. E se discute también en la
Sección~\ref{sec:tradeoffs} como \textit{trade-off} de diseño, ya que su
pregunta de evaluación (¿en qué tamaño de carga compensa la complejidad
$O(\log n)$?) es ortogonal a la elegida y aplica independientemente de qué
extensión funcional se implemente.

\subsection{Diferencias con la implementación sobre Mach 3.0}

El paper original valida \textit{lottery scheduling} integrándolo al
planificador de hilos del \textbf{microkernel Mach 3.0}, reemplazando su
política de prioridades sobre hardware y cargas de trabajo reales: hilos
del kernel, interrupciones de temporizador reales que disparan la
expropiación, bloqueo real en E/S, y múltiples tipos de recurso
(incluyendo memoria y ancho de banda de red) generalizados bajo el mismo
mecanismo de boletos.

Este proyecto es, por diseño, una \textbf{simulación en espacio de
usuario}: no se modifica el \textit{kernel} ni se sustituye el planificador
real de GNU/Linux, que sigue decidiendo los ciclos de CPU físicos por
debajo sin que el programa lo sepa. Las diferencias concretas son:

\begin{itemize}
\item \textbf{Un único CPU lógico simulado}, aplicado mediante el
  protocolo de mutex/condvars de la Sección~\ref{sec:arquitectura}, en vez
  de la expropiación real por interrupción de temporizador del kernel.
\item \textbf{Carga de trabajo sintética e instrumentada} (la serie de
  Taylor de $\pi$) en vez de cargas de programas reales; se eligió
  precisamente porque su estado es verificable, no porque modele algo del
  mundo real.
\item \textbf{Un único tipo de recurso} (unidades de CPU lógicas), sin la
  generalización a memoria o ancho de banda que el paper demuestra.
\item \textbf{Sin cesión real por E/S}: la cesión temprana de la extensión
  se fuerza explícitamente vía \texttt{--yield-config} en vez de surgir de
  un bloqueo de E/S real, porque el prototipo no tiene E/S que modelar.
\end{itemize}

\subsection{Limitaciones del prototipo}

Además de las diferencias estructurales con Mach 3.0, el prototipo tiene
limitaciones propias de su alcance como proyecto de curso: opera sobre un
conjunto fijo de tareas conocido de antemano en el CSV de entrada (el paper
soporta creación dinámica de hilos con boletos heredados); el rango de 5 a
25 tareas es una cota de validación del enunciado, no una limitación
inherente del mecanismo de lotería; la selección de la ganadora es
$O(n)$ sobre la lista de tareas activas (Opción E, no implementada,
discutida como \textit{trade-off} en la Sección~\ref{sec:tradeoffs}); y no
se implementa \textit{ticket transfer} (Opción C), por lo que el prototipo
no demuestra directamente el mecanismo que el paper usa para resolver
inversión de prioridad en llamadas cliente-servidor --- solo su analogía
más cercana, la compensación por cesión temprana.
